% gabarit-poster.tex - Gabarit de poster scientifique A0 (Scriptorium)
%
% Portions adaptees du projet openscience (Synthetic Sciences, InkVell Inc.), Apache-2.0,
% github.com/synthetic-sciences/openscience (backend/cli/skills/writing/latex-posters,
% assets/tikzposter_template.tex). Modifications Marius Yvard, MIT.
%
% Choix de package : tikzposter plutot que beamerposter ou baposter (les trois options source).
% Raisons : classe native autonome (aucune dependance a beamer, contrairement a beamerposter,
% qui detourne un moteur de diaporama pour un objet statique) ; mecanisme \definecolorstyle
% deja pense pour l'injection de couleurs externes (voir bloc CHARTE ci-dessous, calque sur le
% bloc equivalent d'un gabarit-rapport.tex) ; syntaxe de bloc la plus courte des trois (\block{
% Titre}{corps}). Voir produire/references/genre-poster.md pour le playbook complet.
%
% Compilation : xelatex gabarit-poster.tex
%
% Les couleurs de ce gabarit sont remplacees par la sortie de :
%   python3 scripts/theme.py charte-graphique.json --format latex
% (mêmes noms de couleur que gabarit-rapport.tex : ScriptoriumEncre, ScriptoriumTrait,
% ScriptoriumFond, ScriptoriumAccent, ScriptoriumPalUn a ScriptoriumPalQuatre ; le bloc se colle
% integralement, fontspec est charge ci-dessus pour que ses lignes \setmainfont s'appliquent).
%
\documentclass[25pt,a0paper,portrait,margin=10mm,innermargin=15mm,
  blockverticalspace=15mm,colspace=15mm,subcolspace=8mm]{tikzposter}

\usepackage{fontspec}
\usepackage{graphicx}
\usepackage{amsmath}
\usepackage{booktabs}
\usepackage{qrcode}

% ============================================================================
% >>> DEBUT BLOC CHARTE (remplacer par : python3 scripts/theme.py CHARTE.json --format latex)
% ============================================================================
\definecolor{ScriptoriumEncre}{HTML}{2E2A26}
\definecolor{ScriptoriumTrait}{HTML}{8A8175}
\definecolor{ScriptoriumFond}{HTML}{FFFFFF}
\definecolor{ScriptoriumAccent}{HTML}{6E6356}
\definecolor{ScriptoriumPalUn}{HTML}{F4F1EC}
\definecolor{ScriptoriumPalDeux}{HTML}{EEF2F4}
\definecolor{ScriptoriumPalTrois}{HTML}{F1F0EA}
\definecolor{ScriptoriumPalQuatre}{HTML}{F4EEF0}
% ============================================================================
% <<< FIN BLOC CHARTE
% ============================================================================

% Style de couleurs derive de la charte (mecanisme natif tikzposter, bloc de definitions
% intermediaires laisse vide : les couleurs sont deja definies globalement ci-dessus)
\definecolorstyle{ScriptoriumStyle}{
}{
  \colorlet{backgroundcolor}{ScriptoriumPalUn}
  \colorlet{framecolor}{ScriptoriumAccent}
  \colorlet{titlefgcolor}{white}
  \colorlet{titlebgcolor}{ScriptoriumEncre}
  \colorlet{blocktitlebgcolor}{ScriptoriumAccent}
  \colorlet{blocktitlefgcolor}{white}
  \colorlet{blockbodybgcolor}{ScriptoriumFond}
  \colorlet{blockbodyfgcolor}{ScriptoriumEncre}
}
\usetheme{Default}
\usecolorstyle{ScriptoriumStyle}

\title{Titre du poster : concis et descriptif}
\author{Auteur un\textsuperscript{1}, Auteur deux\textsuperscript{2}}
\institute{\textsuperscript{1}Etablissement, ville \quad \textsuperscript{2}Etablissement, ville}

\begin{document}

\maketitle

\begin{columns}

  \column{0.33}
  \block{Introduction}{
    \textbf{Contexte.} Une a deux phrases sur l'importance du sujet.

    \vspace{0.4cm}
    \textbf{Probleme.} Le manque ou la question que le travail adresse.

    \vspace{0.4cm}
    \textbf{Objectif.} Enonce clair de l'objectif du travail.
  }
  \block{Methode}{
    \begin{itemize}
      \item Demarche ou dispositif
      \item Echantillon ou perimetre : n = X
      \item Criteres cles
    \end{itemize}
  }

  \column{0.34}
  \block{Resultat principal}{
    Description breve du resultat le plus important. La figure porte la preuve, le texte
    reste minimal (voir produire/references/genre-poster.md).

    \begin{center}
      % \includegraphics[width=0.9\linewidth]{figure1.pdf}
      \fbox{\parbox{0.85\linewidth}{\centering\vspace{4cm}Emplacement de la figure\vspace{4cm}}}
    \end{center}
    \small Figure 1. Legende avec la donnee cle (n = X, \emph{p} = 0,01).
  }

  \column{0.33}
  \block{Conclusion}{
    \begin{itemize}
      \item Constat un
      \item Constat deux
      \item Portee ou limite
    \end{itemize}
  }
  \block{Pour aller plus loin}{
    \begin{center}
      \qrcode[height=4cm]{https://doi.org/10.0000/exemple}\\
      \small Article, donnees et code
    \end{center}
  }

\end{columns}

\block{}{
  \footnotesize
  \textbf{References.} Reference un ; reference deux.
  \hfill
  \textbf{Contact.} nom@etablissement.fr
}

\end{document}
